\documentclass{article}
\usepackage{amsmath}
\usepackage{amsfonts}
\usepackage{amssymb}
\usepackage{physics}
\usepackage{booktabs}
\usepackage{float}
\usepackage{hyperref}
\usepackage{multirow}
\usepackage{array}
\usepackage[utf8]{inputenc}
\usepackage{graphicx}
\usepackage{caption}
\usepackage{geometry}
\usepackage{xcolor}
\geometry{a4paper, margin=0.8in}
\usepackage{listings}
\definecolor{grey}{rgb}{0.8,0.8,0.8}
\definecolor{darkgreen}{rgb}{0,0.3,0}
\definecolor{darkblue}{rgb}{0,0,0.3}
\def\lstbasicfont{\fontfamily{pcr}\selectfont\footnotesize}
\renewcommand{\lstlistingname}{}
\lstset{%
    numbers=none,
    numberstyle=\small,%
    showstringspaces=false,
    showspaces=false,%
    tabsize=4,%
    frame=lines,%
    basicstyle={\footnotesize\lstbasicfont},%
    keywordstyle=\color{darkblue}\bfseries,%
    identifierstyle=,%
    commentstyle=\color{darkgreen},%\itshape,%
    stringstyle=\color{black},%
    breaklines=true,%
    postbreak=\mbox{\textcolor{red}{$\hookrightarrow$}\space}
}
\lstloadlanguages{C,C++,Java,Matlab,Python,Mathematica}

\title{Comparison of DPC Algorithm Variants}
\author{Xiuyuan Lu}
\date{}

\begin{document}

\maketitle

\section{Introduction}
This document compares the original Density Peak Clustering (DPC) algorithm with three modified versions developed for improved robustness and automation. The comparison covers mathematical formulations, implementation differences, and performance characteristics.

\section{Mathematical Foundations}

\subsection{Original DPC Algorithm}
The original DPC relies on two core metrics:

\begin{itemize}
    \item \textbf{Local Density} ($\rho$):
    \begin{equation}
        \rho_i = \sum_{j \neq i} e^{-\left(\frac{d_{ij}}{d_c}\right)^2}
    \end{equation}
    or alternatively using a cutoff kernel:
    \begin{equation}
        \rho_i = \abs{\{j \mid d_{ij} < d_c\}}
    \end{equation}
    
    \item \textbf{Minimum Distance to Higher Density} ($\delta$):
    \begin{equation}
        \delta_i = \begin{cases}
            \min_{j:\rho_j > \rho_i} d_{ij} & \text{if } \exists j, \rho_j > \rho_i \\
            \max_j d_{ij} & \text{otherwise}
        \end{cases}
    \end{equation}
\end{itemize}

Cluster centers are manually selected from the decision graph ($\gamma_i = \rho_i \times \delta_i$).

\subsection{Modified Versions}
Key improvements in your versions:

\begin{itemize}
    \item \textbf{Density Calculation}:
    \begin{equation}
        \rho_i = \frac{1}{\text{percentile/mean}(\text{KNN distances}) + \epsilon}
    \end{equation}
    
    \item \textbf{Automated Center Selection}:
    \begin{equation}
        n_{\text{clusters}} = \text{KneeLocator}(\gamma \text{ curve})
    \end{equation}
\end{itemize}

\section{Detailed Comparison}

\subsection{Density Calculation}
\begin{table}[H]
\centering
\caption{Density Calculation Methods}
\begin{tabular}{llll}
\toprule
\textbf{Version} & \textbf{Formula} & \textbf{Parameters} & \textbf{Characteristics} \\
\midrule
Original DPC & $\rho_i = \sum e^{-(d_{ij}/d_c)^2}$ & Manual $d_c$ & Sensitive to cutoff \\
Initial & Same as original & Auto $d_c$ (15\% percentile) & Partial automation \\
Modified & $1/\text{KNN percentile}(25\%)$ & $k=2$ neighbors & Local adaptation \\
Final & $1/\text{mean(KNN distances)}$ & $k=5$ neighbors & Smoother but outlier-sensitive \\
\bottomrule
\end{tabular}
\end{table}

\subsection{Cluster Center Selection}
\begin{table}[H]
\centering
\caption{Center Selection Strategies}
\begin{tabular}{lll}
\toprule
\textbf{Version} & \textbf{Method} & \textbf{Automation} \\
\midrule
Original & Manual selection from decision graph & None \\
Initial & Top 8 $\gamma$ points & Fixed count \\
Modified & Top 2 + distance-constrained selection & Semi-automatic \\
Final & KneeLocator + target range (8-9) & Full automation \\
\bottomrule
\end{tabular}
\end{table}

\subsection{Noise Handling}
\begin{table}[H]
\centering
\caption{Noise Processing Approaches}
\begin{tabular}{lll}
\toprule
\textbf{Version} & \textbf{Method} & \textbf{Robustness} \\
\midrule
Original & No explicit handling & Noise affects results \\
Initial & Nearest non-noise assignment & Basic filtering \\
Modified/Final & Density threshold ($\rho > \bar{\rho}$) + nearest neighbor & Improved filtering \\
\bottomrule
\end{tabular}
\end{table}

\section{Performance Characteristics}

\begin{table}[H]
\centering
\caption{Overall Comparison}
\begin{tabular}{llll}
\toprule
\textbf{Feature} & \textbf{Original} & \textbf{Initial} & \textbf{Modified/Final} \\
\midrule
Parameter sensitivity & High (manual $d_c$) & Medium & Low \\
Automation level & None & Partial & High \\
Noise robustness & Poor & Medium & Good \\
High-dimension handling & Poor & Good (with PCA) & Excellent \\
Code complexity & Low & Medium & High/Medium \\
\bottomrule
\end{tabular}
\end{table}

\section{Conclusion}
The evolutionary improvements demonstrate:

\begin{itemize}
    \item Replacement of manual steps with data-driven automation
    \item Enhanced robustness through KNN-based density estimation
    \item Better noise handling via density thresholds
    \item Maintained interpretability while increasing applicability
\end{itemize}

The final version provides the best balance between automation and performance for practical applications.

\section{Appendix}
Here are the source codes for the original version, the modified version, and the final version of the DPC algorithm. 
\lstinputlisting[language=Python, caption={original version}]{../../USdata/clustering_DPC.py}
\lstinputlisting[language=Python, caption={modified version}]{../../USdata/clustering_DPC_modified.py}
\lstinputlisting[language=Python, caption={final version}]{../../USdata/clustering_DPC_remodified.py}

\end{document}